\chapter{Retained contact module}
\label{component:p888d9f79ba3b}
\noindent\textbf{Component }\path{p888d9f79ba3b}\par
{\small\noindent Source: \path{R13/advances/hi/retained_contact_module.tex}\par}
{\small\noindent SHA256: \path{888d9f79ba3b02bd18fe6e06974a16d9ff598ff433f31d3dd360c7d9e4fe2407}\par}
{\small\noindent Inventory use: N0M exact bounded retained remainder, four-monomial circuit, full half-angle conductor and complete normals; only construction imported by full-n theorem\par}
\medskip
\section{Literal \texorpdfstring{$n=0$}{n=0} HI: a retained four-term contact module}
\label{p888d9f79ba3b:sec:r7-hi-n0-retained}
All source and testing variables in this section retain their original
meanings. In particular $q,y,p,v>0$ are independently infinitesimal;
no comparison between a source scale and a testing scale is made.
Write
\[
 d_p=1+p^2,\quad d_v=1+v^2,\quad
 \delta_p=p^2-q,\quad\delta_v=v^2-y,\quad
 S_c=1+R+yt=j+t(1+y)=K+E(1+q).
\]
Use bounded source units
\[
 c=Cyt,\qquad d=Dyj,\qquad \mu=mqK,\qquad a_*=\alpha qE.
\]
At literal $n=0$, the exact source formulas give
\begin{align}
 2AH/(qy)&=R, & \bar R&=1,\\
 c-d&=\frac y2\{R(1+q)-t\},&
 \mu-a_*&=\frac q2\{E-R(1+y)\}.                 \tag{N0M.1}
\end{align}
The four starred weights have residue $1/4$, even when any positive
opening tends to zero. The last two identities follow from the rational
source formulas; before imposing $n=0$ their respective right sides are
$y\{(2n+1)R(1+q)-t\}/2$ and
$q\{E-(2n+1)R(1+y)\}/2$.

\paragraph{A common state with exact contact.}
Set
\[
 A_0=W-pY,\qquad B_0=vY-Z,\qquad
 S=X-vW-pZ+pvY,\qquad s=-S/S_c,
 \qquad Q=\frac{\sqrt{qy}}H Y.
\]
Then $J=d_vA_0+vS$ and $M=d_pB_0-pS$. Define
\[
 N_c=\frac{A_0-tv s}{\sqrt y},\qquad
 N_m=\frac{B_0+Ep s}{\sqrt q},\qquad
 P=\frac p{\sqrt q},\quad V=\frac v{\sqrt y},\quad T=PV,
\]
and the additional contact normals
\[
 N_D=d_vN_c+t\delta_vVs,\qquad
 N_M=d_pN_m-E\delta_pPs.                         \tag{N0M.2}
\]
At the actual source contact, $P=V=1$, $s=Q=\sqrt q$, and
$N_c=N_m=N_D=N_M=0$. All definitions are tensor identities; neither
$P$ nor $V$ is assumed bounded.

\paragraph{Retain all singular opening excesses.}
The following four positive squares are retained in the parent:
\begin{align}
 \mathcal E={}&\frac{4cj d_v}{tS_c}N_c^2
 +\frac{4dt(1+y)}{jS_c}N_D^2\notag\\
 &+\frac{4\mu E(1+q)}{KS_c}N_M^2
 +\frac{4a_*K d_p}{ES_c}N_m^2.                 \tag{N0M.3}
\end{align}
Their coefficients are products of positive pure-source and pure-testing
factors, so the displayed terms are actual tensor squares. They need
not be bounded. No excess is discarded before collecting the remainder.

Put $g=S_c-tv^2-Ep^2$. The unchanged full parent has the exact identity
\begin{align}
 F-\mathcal E={}&h_cN_c^2+h_mN_m^2
 -4\sqrt{qy}\,pv d_pd_vN_cN_m\notag\\
 &+N_c\{C_sVs+4H\sqrt y\,v d_pd_vTQ\}
 +N_m\{M_sPs-4H\sqrt q\,p d_pd_vTQ\}
 +\mathcal G,                                           \tag{N0M.4}
\end{align}
where
\begin{align*}
 h_c&=\frac{4d_v}{S_c}\{c(1+y)+dd_v\}
       +2yd_v(1+v^2+p^2v^2),\\
 h_m&=\frac{4d_p}{S_c}\{\mu d_p+a_*(1+q)\}
       +2qd_p(1+p^2+p^2v^2),\\
 C_s&=4d_v\left[y\{R(1+q)-d_pg\}
                    +\frac{2dt}{S_c}\delta_v\right],\\
 M_s&=4d_p\left[q\{d_vg-R(1+y)\}
                    -\frac{2\mu E}{S_c}\delta_p\right].
\end{align*}
The two diagonal coefficients have residue one. The normal--normal
off-diagonal coefficient and all four raw coefficients in the first
normal line are infinitesimal bounded tensors. In particular no inverse
opening occurs in (N0M.4).

The scalar graph is
\[
 \mathcal G=f_AP^2s^2+f_BV^2s^2
       +f_CT^2Q^2+f_DQ^2-eTsQ,                         \tag{N0M.5}
\]
with
\begin{align*}
 f_A={}&4\left[\mu K+a_*Ed_p-2\mu E\delta_p
                    +\frac{\mu E^2}{S_c}\delta_p^2\right]
 +2qd_v(E\delta_p-K+tv^2)^2-2E^2qv^2d_p,\\
 f_B={}&4\left[ctd_v+dj-2dt\delta_v
                    +\frac{dt^2}{S_c}\delta_v^2\right]
 +2y(t\delta_v-j)^2,\\
 f_C={}&2H^2d_pd_v,\qquad f_D=2HRd_pd_v,
 \qquad e=4Hd_pd_vg.
\end{align*}
All five are polynomials in $p,v$ with bounded source coefficients.
Their residues are, respectively, $2,2,2,2,8$. This assertion is uniform
over every opening permitted by the exact HI source entry.

\paragraph{The exact source contact jets.}
A star denotes $p=\sqrt q$, $v=\sqrt y$. Set
$\sigma=f_A+f_B+f_C+f_D-e$. The rational source formulas give
\begin{align}
 f_A^*+f_B^*&=f_C^*+f_D^*=e^*/2,\notag\\
 q^{1/4}(\partial_{\sqrt p}\sigma)^*
   &=-4(f_A^*-f_D^*),\notag\\
 y^{1/4}(\partial_{\sqrt v}\sigma)^*
   &=-4(f_B^*-f_D^*).                              \tag{N0M.6}
\end{align}
Also
\[
 C_s^*+4Hy(1+q)(1+y)=0,\qquad
 M_s^*-4Hq(1+q)(1+y)=0.                           \tag{N0M.7}
\]
These are actual-contact identities, not residue identities. The checker
verifies the only source-dependent ones modulo the exact numerator of
$n$, rather than replacing $R$ by its residue.

\paragraph{Finite contact lift of the graph and both state normals.}
Here the SOS is a tensor SOS at each original pure-testing state. Its
roots need not be linear polynomials in $W,X,Y,Z$. If $sQ>0$, a common
sign change permits $s,Q>0$. This is a pure-testing sign choice, because
$sQ$ is a positive source factor times $-SY$. Put
\[
 A=Ps,\quad B=Vs,\quad C=TQ,\quad D=Q,
 \quad G=(ABCD)^{1/4}=\sqrt{TsQ},\quad E_i=A_i-G.
\]
Every $A_i$ and $G$ is a positive decomposable tensor. Thus every radical
below is a product of a pure-source radical and a pure-testing radical.
The basic circuit is
\[
 L=A^2+B^2+C^2+D^2-4G^2
   =\sum_iE_i^2+2G\sum_iE_i.                     \tag{N0M.8}
\]
For each pairing $ij|k\ell$,
\[
 G\sum_iE_i=G(\sqrt{A_i}-\sqrt{A_j})^2
 +G(\sqrt{A_k}-\sqrt{A_\ell})^2
 +2G\{(A_iA_j)^{1/4}-(A_kA_\ell)^{1/4}\}^2.
\]
Moreover,
\begin{align*}
 L&=(A-C)^2+(B-D)^2+2(\sqrt{AC}-\sqrt{BD})^2,\\
 (r-r')^2&=(r-G)^2+(r'-G)^2
             +2G(\sqrt r-\sqrt{r'})^2
 \quad(r=\sqrt{AC},\ r'=\sqrt{BD}),
\end{align*}
and the analogous identities hold for $BC|AD$. A finite average of
these decompositions gives strictly positive constant diagonal pivots
simultaneously on the $E_i$, the variance roots in (N0M.8), and
\[
 F_p=\sqrt{AC}-G,\qquad F_v=\sqrt{BC}-G.
\]
The exact half-angle conductors are
\[
 G(\sqrt p-q^{1/4})=q^{1/4}F_p,
 \qquad G(\sqrt v-y^{1/4})=y^{1/4}F_v.          \tag{N0M.9}
\]
Their coefficients are pure-source infinitesimals.

For completeness, the first-jet cancellation is explicit. Choose the
positive source unit $k=(f_C^*+f_D^*)/2\sim2$. Taylor-expand the bounded
coefficient polynomials in $z=\sqrt p,w=\sqrt v$ around
$z_0=q^{1/4},w_0=y^{1/4}$. All divided-difference coefficients remain
bounded, since these are monic polynomial divisions by pure-source
centers. Explicitly each monomial has the finite identity
\[
 z^iw^j=\sum_{a=0}^i\sum_{b=0}^j
 {i\choose a}{j\choose b}z_0^{i-a}w_0^{j-b}
 (z-z_0)^a(w-w_0)^b.
\]
Collecting total degrees zero, one, and at least two is the required
exact Taylor certificate. Its coefficients are bounded because all
exponents are nonnegative and all source coefficient denominators are
positive source units. No analytic remainder or mixed division enters.
After (N0M.9), every term in $\mathcal G-kL$ is an infinitesimal
bounded quadratic in the listed roots, except possibly
\[
 G\left[2\sum_i(f_i^*-k)E_i
       +z_0\sigma_z^*F_p+w_0\sigma_w^*F_v\right].
\]
By (N0M.6) this equals
\begin{align*}
 2(f_D^*-k)G\sum_iE_i
 &-\frac{z_0\sigma_z^*}{2}
       \{\sqrt G(\sqrt A-\sqrt C)\}^2\\
 &-\frac{w_0\sigma_w^*}{2}
       \{\sqrt G(\sqrt B-\sqrt C)\}^2.
\end{align*}
All three coefficients are source infinitesimals, and all their roots
already have strict pivots in the finite averaged circuit.

Finally the $N_c$ coupling in (N0M.4) is
$N_c(C_sB+\beta_cC)$, where
$\beta_c=4H\sqrt y\,v d_pd_v$ is infinitesimal.
The sum $C_s+\beta_c$ vanishes at the source contact by (N0M.7).
Writing $B=G+E_B$, $C=G+E_C$, monic-dividing that coefficient sum in
$z-z_0,w-w_0$, and applying (N0M.9) expresses this entire coupling by
infinitesimal bounded entries on $N_c,E_B,E_C,F_p,F_v$. The $N_m$
coupling is identical, using $A,C$ and
$\beta_m=-4H\sqrt q\,p d_pd_v$. Thus the two state normals have an
identity-residue block, and the graph has the strict circuit block.
The established strict bounded-Gram absorption lemma now gives an
actual tensor SOS of $F-\mathcal E$.

If $sQ<0$, the graph is instead a strict bounded diagonal form on the
four raw roots and $\sqrt{T|sQ|}$, with residue coefficients
$2,2,2,2,8$. The first-normal couplings already have infinitesimal raw
coefficients. If $sQ=0$, simply omit the last root. The same bounded-Gram
lemma applies directly. Consequently this argument covers literal
$n=0$ throughout the critical HI testing corner, at arbitrary independent
source/testing rates and through all positive source openings. It does
not assert the corresponding $n>0$ extension.

\paragraph{Verification scope.}
\path{check_retained_contact_module.py} checks the pair excess
identities, the complete signed core, the four-term graph, the exact
source differences, and all contact-jet identities. Its 18 checks are
exact symbolic equalities. The finite-radical lift and the uniform
boundedness argument are the analytic argument above, not conclusions
drawn from numerical sampling.
